\documentclass[12pt]{article}
\usepackage[margin=1in]{geometry}
\usepackage{setspace}
\usepackage{amsmath, amssymb, amsthm}
\usepackage{titlesec}
\usepackage{enumitem}

\title{\textbf{Reviewer Report}\\ \Large Political Polarization and Market Reactions to Auditor-Change Disclosures}
\author{Reviewer}
\date{\today}

\begin{document}
\maketitle
\doublespacing

\section*{General Assessment}
This paper examines how state-level political polarization impacts market reactions to auditor-change disclosures (Form 8-K Item 4.01). The premise---that institutional trust, mediated by political polarization, generates heterogeneous priors and thereby amplifies trading volume and price movements---is a creative and novel integration of political economy and accounting research. 

However, the manuscript suffers from a critical disconnect between the primary theoretical predictions and the empirical findings, particularly regarding the ambiguity mechanism. Furthermore, the theoretical model, while mathematically sound, is somewhat elementary and could be streamlined. The prose contains repetitive defensive phrasing that detracts from the paper's flow. 

To provide a rigorous evaluation, I have structured this review in five iterative passes, moving from macro-level contributions down to micro-level mathematical and notational consistency.

% ── Iteration 1 ───────────────────────────────────────────────────────────────

\section{Iteration 1: Novelty and Contribution Assessment}

The integration of the identity/alienation literature from political science into the auditing literature is the paper's strongest asset. Showing that mandatory disclosure effectiveness is bounded by the political-information environment is a compelling contribution to the disclosure literature.

\textbf{Deficiencies \& Actionable Suggestions:}
\begin{itemize}
    \item \textbf{The Identification Strategy:} The mechanism strictly relies on the assumption that the marginal trader evaluating these 8-Ks is a local retail investor subject to local state political polarization. However, the paper uses 678 medium-to-large firms covered by CRSP/Compustat. Institutional investors, who are geographically dispersed, dominate trading in these markets. 
    \item \textbf{Action:} You must condition the main effect on institutional ownership. Merge your sample with 13F filings. If your theory holds, the polarization coefficient ($\beta$) must be significantly larger (or exclusively present) in the subset of firms with high retail ownership / low institutional ownership. 
\end{itemize}

% ── Iteration 2 ───────────────────────────────────────────────────────────────

\section{Iteration 2: Theory-to-Empirics Gaps}

There is a significant conceptual gap between the theoretical construct of ``polarization'' and the empirical proxy.

\textbf{Deficiencies \& Actionable Suggestions:}
\begin{itemize}
    \item \textbf{Proxy Mismatch:} The theoretical model relies on a mean-preserving spread of priors (a widening of $|\pi_D - \pi_R|$). The empirical proxy is $1 - |D-R|$ (state-level electoral competitiveness). A 50-50 state means the local population is evenly split. It does \textit{not} necessarily mean that the specific investors holding the firm's stock are polarized, nor does it guarantee that the distance between the two ideological camps is wider than in a 70-30 state.
    \item \textbf{Action:} Acknowledge this explicitly in the empirics section, not just as ``noise.'' Better yet, utilize the DW-NOMINATE cross-sectional variance or county-level variance (which you relegated to a robustness check) as a co-primary measure, as these better capture actual ideological distance rather than just competitive vote shares.
\end{itemize}

% ── Iteration 3 ───────────────────────────────────────────────────────────────

\section{Iteration 3: Logical Coherence \& The Ambiguity Contradiction}

This is the most critical flaw in the current draft. The paper establishes abnormal volume ($AbVol$) as the primary formal prediction of the disagreement model (P1). It then hypothesizes that ambiguity amplifies this effect (P3). 

\textbf{Deficiencies \& Actionable Suggestions:}
\begin{itemize}
    \item \textbf{The P3 Failure:} In Section 6.4, the empirics show that volume is driven by the \textit{low-ambiguity} subsample, failing the P3 test. The author writes: \textit{``One possible interpretation is that low-ambiguity events... give investors with opposing political priors concrete grounds on which to disagree... This is a post-hoc rationalization, however...''} 
    \item \textbf{Critique:} Candidness is appreciated, but a post-hoc rationalization that explicitly contradicts the theoretical framework is fatal in top-tier review. If unambiguous signals generate more volume, then your theory in Section 3.4 (which claims ambiguity shifts the locus of disagreement to subjective priors) is either misspecified or empirically invalid.
    \item \textbf{Action:} You must resolve this contradiction before submission. Consider revising the theoretical framework to incorporate \textit{directional} bias alongside precision. For example, a clear negative signal (low ambiguity) might trigger differential disposition effects among polarized investors. Alternatively, rethink your empirical definition of \texttt{high\_ambiguity}. Is a lateral Big-4-to-Big-4 switch genuinely ``ambiguous,'' or just uninformative noise?
\end{itemize}

% ── Iteration 4 ───────────────────────────────────────────────────────────────

\section{Iteration 4: Mathematical Verification and Notation}

I have carefully audited the proofs in Appendix A. 

\textbf{Deficiencies \& Actionable Suggestions:}
\begin{itemize}
    \item \textbf{Proof Validation:} The proofs for Propositions 1, A1, A2, and A3 are algebraically correct. The variance of the two-point distribution $\mathrm{Var}_i[\phi(\pi_i)] = \omega_D \omega_R \bigl(\phi(\pi_D) - \phi(\pi_R)\bigr)^2$ is derived cleanly.
    \item \textbf{Critique of Theoretical Depth:} The math, while correct, is trivial. Propositions A1, A2, and A3 are essentially the same mathematical statement repeated three times. They do not need to be separate propositions.
    \item \textbf{Action:} Condense Propositions A1, A2, and A3 into a single unified proposition (similar to your current Prop A3) in the main text. State that the latent index $X_g$ captures priors, trust, and perceived precision, and prove the variance property once. This will make the theory section read as a more cohesive structural model rather than a step-by-step textbook derivation.
    \item \textbf{Notation Consistency:} In Equation 6, you use $X_{e,t(e)-1}$ for firm controls, but in Section 3.4 you use $X_g$ for the latent log-odds index. Overloading $X$ causes confusion. Change the firm controls matrix in Eq 6 to $C_{e, t(e)-1}$ or $Z_{e, t(e)-1}$.
\end{itemize}

% ── Iteration 5 ───────────────────────────────────────────────────────────────

\section{Iteration 5: Redundancy, Flow, and Writing}

The manuscript frequently repeats defensive arguments, which dilutes the impact of the findings.

\textbf{Deficiencies \& Actionable Suggestions:}
\begin{itemize}
    \item \textbf{Repetitive Phrasing:} The phrase ``generic attention or salience stories'' (or close variants) appears repeatedly in the Introduction, Section 3.4, Section 6.3, Section 6.4, and the Conclusion. Similarly, the phrase ``objective informational guidance'' is heavily overused.
    \item \textbf{Action:} Streamline the text. Make the argument against the ``salience'' alternative once, forcefully, when introducing the placebo test and ambiguity splits, and then trust the reader to remember it. 
    \item \textbf{Sample Size:} $N=678$ is quite small for a 22-year window (2001-2023). Briefly clarify in the data section whether this is due to the exclusion of 8-K amendments, or missing CRSP/Compustat matches, and explicitly defend the statistical power of your tests given the $N$.
\end{itemize}

\end{document}